\chapter{Semantic Mapping Rules}
\label{appendix:semantic_labeling_rules}

This appendix documents the rule-based semantic mapping used to assign interpretable labels to the latent states of the selected S2A4 hierarchical model.

\section{Threshold Values}

\begin{table}[htbp]
\centering
\small
\caption{Threshold values used by the semantic mapping rules.}
\label{tab:semantic_mapping_thresholds}
\begin{tabular}{l l}
\toprule
Parameter & Value \\
\midrule
\texttt{lc\_flag\_thresh} & 0.5 \\
\texttt{lateral\_velocity\_threshold} & 0.1 \\
\texttt{lateral\_acceleration\_threshold} & 0.1 \\
\texttt{lateral\_velocity\_slope\_threshold} & 0.003 \\
\texttt{ay\_zero\_frac\_threshold} & 0.5 \\
\texttt{brake\_ax\_neg\_frac} & 0.90 \\
\texttt{brake\_ax\_last} & -0.15 \\
\texttt{tactical\_brake\_ax\_last} & -0.25 \\
\texttt{tactical\_brake\_vx\_slope} & -0.01 \\
\texttt{tactical\_brake\_thw} & 2.0 \\
\texttt{tactical\_brake\_ttc} & 70.0 \\
\texttt{interaction\_thw\_tight} & 2.0 \\
\texttt{interaction\_ttc\_tight} & 60.0 \\
\texttt{strong\_acc\_front\_exists\_min} & 0.50 \\
\texttt{strong\_acc\_ax\_pos\_frac} & 0.90 \\
\texttt{strong\_acc\_ax\_last} & 0.15 \\
\texttt{constrained\_acc\_ax\_pos\_frac} & 0.65 \\
\texttt{constrained\_acc\_ax\_last} & 0.08 \\
\texttt{constrained\_acc\_front\_exists} & 0.80 \\
\texttt{free\_flow\_front\_exists\_max} & 0.50 \\
\texttt{constrained\_follow\_front\_exists} & 0.80 \\
\texttt{constrained\_follow\_vx\_last\_max} & 21.0 \\
\texttt{constrained\_follow\_jerk\_min} & 0.75 \\
\texttt{stable\_follow\_front\_exists} & 0.80 \\
\texttt{stable\_follow\_vx\_slope\_abs\_max} & 0.008 \\
\texttt{stable\_follow\_ax\_last\_abs\_max} & 0.20 \\
\bottomrule
\end{tabular}
\end{table}


\section{Priority Order}

When more than one rule condition is satisfied for the same window, the final semantic label is assigned according to the following precedence order:
\begin{enumerate}
    \item lane change,
    \item braking, followed by the split into tactical brake versus mild brake,
    \item strong acceleration,
    \item constrained acceleration,
    \item constrained following,
    \item stable following,
    \item free flow,
    \item unknown.
\end{enumerate}
This priority order ensures that rarer or more distinctive maneuvers, such as lane changes and braking events, are not overridden by broader longitudinal behavior rules.
 
\section{Composite Definitions}
\par The condition \texttt{tactical\_brake} is true if at least one strong braking cue is present:\\
\texttt{ax\_last <= tactical\_brake\_ax\_last},\\
\texttt{vx\_slope <= tactical\_brake\_vx\_slope},\\
\texttt{front\_thw\_last <= tactical\_brake\_thw} when the feature is available, or\\
\texttt{front\_ttc\_min <= tactical\_brake\_ttc} when the feature is available.\\

The condition \texttt{tight\_interaction} is true if\\
\texttt{front\_exists\_frac >= stable\_follow\_front\_exists}\\
and at least one interaction-risk cue is active:\\
\texttt{front\_thw\_last <= interaction\_thw\_tight} when available, or\\
\texttt{front\_ttc\_min <= interaction\_ttc\_tight} when available.\\

\subsection*{Lane-change rule}

The lane-change rule has the highest priority. 
A window is labeled as \texttt{lane\_change} if either the left or right lane-change indicator exceeds the threshold. 
It can also be triggered by a sustained lateral-motion pattern defined through lateral velocity, lateral acceleration, lateral velocity slope, and the fraction of near-zero lateral acceleration within the window.

\subsection*{Braking rules}

Braking is detected when negative longitudinal acceleration dominates most of the window, the final longitudinal acceleration is sufficiently negative, and a front vehicle is present for a large fraction of the window. 
Once braking is detected, the window is split into either \texttt{mild\_brake} or \texttt{tactical\_brake}. 
If any tactical braking cue is present, the label becomes \texttt{tactical\_brake}; otherwise it is assigned as \texttt{mild\_brake}.

\subsection*{Acceleration rules}

Two acceleration-related labels are defined. 
A window is labeled \texttt{strong\_acceleration} when positive longitudinal acceleration dominates the window, the final acceleration is clearly positive, and either a front vehicle is present or the situation is already classified as a tight interaction. 
A window is labeled \texttt{constrained\_acceleration} when positive acceleration is present under stronger leader presence, indicating that the maneuver is still shaped by traffic context.

\subsection*{Following rules}

Two following-related labels are defined. 
A window is labeled \texttt{constrained\_following} when a front vehicle is present for most of the window and either the motion shows low-speed, high-jerk behavior or the interaction is tight. 
This rule is intended to capture pressured or corrective following behavior. 
A window is labeled \texttt{stable\_following} when a front vehicle is present, the speed slope and longitudinal acceleration remain close to zero, and the interaction is not tight. 
This rule therefore represents smoother and more stable car-following behavior.

\subsection*{Free-flow rule}

If none of the higher-priority maneuver rules fires and the fraction of front-vehicle presence remains below the free-flow threshold, the window is labeled \texttt{free\_flow}. 
This captures low-interaction motion with little direct leader influence.

\subsection*{Fallback case}

If none of the semantic rules is satisfied, the window is assigned the fallback label \texttt{unknown}. 
This case corresponds to windows for which no lane-change, braking, acceleration, or following pattern is detected by the rule set.